\documentclass[11pt]{article}
\usepackage[margin=1in]{geometry}
\usepackage{amsmath,amssymb,amsthm,hyperref,enumitem}
\hypersetup{colorlinks=true,urlcolor=blue,linkcolor=blue}
\newtheorem{theorem}{Theorem}
\newtheorem{lemma}[theorem]{Lemma}
\title{Monotonicity of the signature of a mixed graph matrix}
\author{Principia Math In Progress}
\date{22 September 2026}
\begin{document}
\maketitle
\noindent\textbf{MathDB 364187. Task 1: a full proof.}
We prove Conjecture 1 of Bonetto and Jard\'on Kojakhmetov
\cite[p.~7]{source}, also retained as Conjecture 3.1.1 in
\cite[p.~55]{thesis}. The proof applies to the more general symmetric
matrix family proposed immediately after their conjecture.

For a real symmetric matrix $Q$, let $n_+(Q)$, $n_0(Q)$ and $n_-(Q)$
be its positive, zero and negative eigenvalue counts, with multiplicity.
Its signature is $s(Q)=n_+(Q)-n_-(Q)$. Write $X\preceq Y$ if
$Y-X$ is positive semidefinite.

\begin{theorem}\label{main}
Let $A,B$ be real symmetric $n\times n$ matrices, with $A\succeq0$.
For $M(t)=A+tB$ and $0<u<v$,
\[
 n_+(M(v))\le n_+(M(u)),\qquad
 n_-(M(v))\ge n_-(M(u)).                              \tag{1}
\]
Consequently $s(M(v))\le s(M(u))$, including at singular parameter values.
\end{theorem}

\begin{lemma}\label{order}
If $X,Y$ are real symmetric matrices and $X\preceq Y$, then
$n_+(X)\le n_+(Y)$ and $n_-(X)\ge n_-(Y)$.
\end{lemma}
\begin{proof}
The number $n_+(Q)$ is the largest dimension of a linear subspace on
which $x^TQx$ is strictly positive for every nonzero $x$. Indeed, the
positive eigenspace has that dimension. Any subspace of larger dimension
intersects the span of the nonpositive eigenvectors nontrivially, so
cannot be positive definite. The analogous statement holds for $n_-(Q)$.

If $x^TXx>0$ on a subspace away from zero, then
$x^TYx\ge x^TXx>0$ there. This proves the assertion for $n_+$.
If $x^TYx<0$ on a subspace away from zero, then
$x^TXx\le x^TYx<0$ there, proving the assertion for $n_-$.
Zero eigenvalues cause no change to either argument.
\end{proof}

\begin{proof}[Proof of Theorem~\ref{main}]
Positive scalar multiplication preserves all three inertia counts.
Normalize the matrix by its parameter:
\[
 N(t)=\frac{M(t)}t=B+\frac At.
\]
For $0<u<v$ we have
\[
 N(u)-N(v)=\left(\frac1u-\frac1v\right)A\succeq0.     \tag{2}
\]
Thus $N(v)\preceq N(u)$. Lemma~\ref{order} gives (1), since
$N(t)$ and $M(t)$ have the same inertia. Subtracting the two count
inequalities proves the signature inequality.
\end{proof}

\section*{Application to the source conjecture}
In Definition 3 of \cite{source}, a simple undirected graph $F$ has
its edges partitioned between spanning subgraphs $G$ and $H$, and
\[
 M(\epsilon)=L_G+\epsilon A_H,\qquad \epsilon>0.
\]
Here $A_H$ is real symmetric, and the graph Laplacian satisfies
\[
 x^TL_Gx=\sum_{\{i,j\}\in E_G}(x_i-x_j)^2\ge0.
\]
Theorem~\ref{main} with $A=L_G$ and $B=A_H$ therefore proves the
conjecture for every permitted graph and every positive parameter.
In the associated gradient system $\dot q=-M(\epsilon)q$, the unstable
eigenspace has dimension $n_-(M(\epsilon))$, which is nondecreasing.
The stable dimension $n_+(M(\epsilon))$ is nonincreasing.

As in the source's step-function plots, decreasing means nonincreasing;
constant intervals are allowed. The theorem includes the exact values at
zero crossings and allows a common kernel for every parameter. It does
not assert that each unnormalized eigenvalue is a decreasing real-valued
function. The positive-parameter domain is precisely that of the source.

\section*{Verification and attribution}
Lemma~\ref{order} is the standard order property of inertia, with an
elementary proof included. The only further step is the normalization (2).
The exact verifier checks all graph partitions through four vertices,
general symmetric pencils, and singular and repeated crossings with a
permanent kernel. Independent characteristic-root isolation checks the
inertia routine. These finite tests supplement the unrestricted proof.
Review is by the authoring agent only.

\section*{Proposed follow-up Tasks 2--5}
\begin{enumerate}[label=\textbf{Task \arabic*.},start=2,leftmargin=*]
\item Determine the positive crossing parameters and their multiplicities
from a canonical decomposition of the pencil $(A,B)$, allowing a
nontrivial common kernel.
\item Characterize graph partitions $(G,H)$ attaining each possible total
increase of $n_-(L_G+\epsilon A_H)$ on $(0,\infty)$.
\item Extend the inertia monotonicity statement to families
$a(t)A+b(t)B$ with $b(t)>0$, and characterize the coefficient functions
that guarantee it for every $A\succeq0$ and symmetric $B$.
\item Determine hypotheses for an analogous monotonicity theorem for
self-adjoint operator pencils, including domains and finiteness of
negative spectral subspaces.
\end{enumerate}

\begin{thebibliography}{9}
\raggedright
\bibitem{source} R. Bonetto and H. Jard\'on Kojakhmetov,
\emph{On the Eigenvalues of Graphs with Mixed Algebraic Structure},
\href{https://arxiv.org/abs/2408.00487}{arXiv:2408.00487v1}, 2024,
Definition 3, Definition 8 and Conjecture 1.
\bibitem{thesis} R. Bonetto,
\emph{Oscillations and diffusion in dynamic networks models},
University of Groningen PhD thesis, 2025,
\href{https://doi.org/10.33612/diss.1329052104}{doi:10.33612/diss.1329052104},
Section 3.1, pp. 47--58, especially Conjecture 3.1.1 on p. 55.
\end{thebibliography}
\end{document}
